\chapter{Introdução}

A histopatologia é responsável pelo diagnóstico precoce de patologias de tecidos agressivas como cânceres e doenças autoimunes, infecciosas e degenerativas. Geralmente esses diagnósticos são feitos manualmente por patologistas e demandam um processo de análise de extração de tecidos biológicos, coloração por meio de reagentes químicos (para realçar o contraste de estruturas celulares) e então a investigação microscópica de especialistas, e todo esse processo é custoso e pode ser acelerado por computadorização \cite{autoHistImg}. Além do processo custoso, a falta de padronização entre laboratórios, a requisição de uma exaustiva anotação de imagens de tecidos, junto de restrições éticas e de privacidade de dados, dificultam a criação de grandes bases de imagens para utilização em modelos de aprendizado profundo \cite{histImgColorectal}. Apesar da escassez de imagens, técnicas como transferência de aprendizado, aumento de dados e extração de \textit{features} estão sendo utilizadas na literatura para tornar possível a aplicação de modelos de \textit{deep learning}, e os resultados obtidos mostram que em determinados casos os modelos superaram a precisão no diagnóstico quando comparados a especialistas da área \cite{bejnordi2017diagnostic, esteva2017dermatologist}.

As arquiteturas baseadas em aprendizado profundo são aplicadas para mitigar as limitações impostas pela avaliação humana. As Redes Neurais Convolucionais (CNNs) tornaram-se predominantes na automatização da extração de características locais em imagens médicas. A análise de lâminas inteiras exige frequentemente a divisão em pequenas áreas independentes, uma abordagem descontextualizada que viola o comportamento diagnóstico natural do patologista e compromete a perceção global do tecido \cite{histImgColorectal}. Para capturar correlações de longa distância, os \textit{Vision Transformers} (ViTs) foram introduzidos na patologia digital \cite{shamshad2023transformers}. A combinação das capacidades das CNNs e dos ViTs resultou no desenvolvimento de modelos híbridos. Os resultados indicam que a união da extração de características locais com a atenção global proporcionada pelos modelos híbridos produziu um desempenho superior na análise de imagens histopatológicas quando comparado ao uso isolado de redes convolucionais ou transformadores \cite{oliveira2025hybrid, dai2021coatnet}.

A modelagem topológica oferece uma via alternativa para a extração de \textit{features}. Fundamentada na teoria dos grafos \cite{euler1736solutio}, a abstração de redes complexas mapeia a dinâmica de sistemas intrincados \cite{BOCCALETTI2006175}. Ao modelar os pixels de uma imagem como vértices e estabelecer arestas baseadas em distâncias métricas, descritores topológicos são obtidos para auxiliar na classificação de texturas \cite{BACKES2013168, CNT}. O cálculo do peso dessas arestas a partir de diferenças de intensidade aprimorou a obtenção das características \cite{ribas2020fusion, ribas2024colorspectral}. A expansão progressiva do raio de conexão entre os nós gera uma sequência ordenada de descritores, caracterizando uma dinâmica na qual as macro-organizações dependem da distribuição dos componentes microscópicos. Para processar essa evolução histórica, as Redes Neurais Recorrentes (RNNs) interpretam a variação do grafo como passos sequenciais e geram representações discriminatórias focadas na organização do tecido \cite{hochreiter1997long}.

A integração de múltiplas representações de dados eleva a precisão analítica de sistemas de inteligência artificial. Arquiteturas de dupla via (\textit{Two-Stream Networks}) processam diferentes fontes de informação em paralelo para extrair representações distintas \cite{simonyan2014two, yousaf2024twostream}. A combinação de características visuais com assinaturas topológicas é executada de maneira eficiente por meio da fusão tardia (\textit{Late Fusion}). A estratégia de fusão tardia atua diretamente nas camadas de decisão dos modelos e permite que as representações independentes sejam combinadas sem interferir na estabilidade do treinamento inicial \cite{kittler1998combining, wang2023stnet}. Na patologia computacional, a concatenação das características baseadas em conhecimento de domínio com atributos profundos tem provado ser uma solução eficaz para o treinamento de modelos precisos sob condições de escassez de dados \cite{baltrusaitis2019multimodal}.

A literatura carece de abordagens que combinem representações visuais extraídas por modelos híbridos com assinaturas topológicas sequenciais em imagens médicas. A aplicação direta de métodos clássicos de extração de textura ignora as relações entre os canais de cor presentes nos exames histopatológicos originais. A motivação deste trabalho reside na hipótese de que a união da representação visual gerada por um modelo híbrido com a evolução topológica de uma rede complexa processada por uma arquitetura recorrente resulta em um classificador de alta acurácia. A justificativa para essa integração pauta-se na necessidade de contornar a escassez de imagens anotadas através do enriquecimento das características extraídas em cada amostra disponível.

As contribuições deste trabalho incluem a adaptação da modelagem de redes complexas para imagens histopatológicas coloridas e o desenvolvimento de uma arquitetura de dupla via baseada em fusão tardia. O modelo proposto processa as características morfológicas por meio de uma rede híbrida baseada em convolução e atenção em paralelo com o processamento da dinâmica topológica do tecido utilizando uma arquitetura recorrente. O sistema desenvolvido foi avaliado em bases de dados histopatológicas públicas. Os resultados demonstram a eficácia da abordagem proposta na classificação de tecidos celulares frente a modelos de referência consolidados.

% REFORMULACAO DA INTRODUCAO

% 1° paragrafo (geralzao): histopatologia (definicao, importancia para identificacao de doencas) -> uso de IA pra classificacao de tecidos -> dificuldades desse contexto, com comentarios sobre a escassez de imagens médicas no geral -> extração de features -> ação conjunta de redes de duas vias para aprimoramento dessas atividades

% 2° paragrafo (revisao sobre aplicacao de modelos em histopatologia) | falar de cnn, vit e hibridos (refs com superioridade de hibridos) | 

% 3° paragrafo (revisao sobre extracao de features) | utilizacao de extracao de textura em histopatologia e geral tb

% 4° paragrafo (revisao sobre fusao tardia) | mostara modelos de fusao tardia apicados na histopatologia e em outras areas | 

% 5° paragrafo (gap, motivacao e justificativa)

% 6° contribuicoes
%
% \subsection{antiga introdução}

% deep learning | escassez de imagens médicas | data aug como forma de mitigar overfitting e escassez
% essa introducao explica genericamente o avanco do deep learning, falta dizer como anda fazendo, em quais dominios medicos e como (diagnostico, segmentacao, oq???)
% a escassez eh ate ok, porem nao diz imagens do que, como sao feitas a criacao das iamgens, falta direcionamento a histopatologia
% o data aug tende a dizer que irei propor/comparar formas de data aug, mas na verdade o foco aqui é extracao de features, nao data aug...
% O avanço das arquiteturas de aprendizado profundo revolucionou a análise computacional ao permitir a extração de características semânticas complexas com alta precisão. Contudo, no domínio da saúde, a aplicação dessas técnicas frequentemente esbarra na escassez de imagens médicas anotadas, decorrente de rigorosos protocolos de privacidade, altos custos e da necessidade de especialistas para rotulação \cite{litjens2017survey}. Para mitigar essa limitação, o aumento de dados (\textit{Data Augmentation}) consolidou-se como uma etapa essencial. Ao expandir o espaço amostral, essas técnicas atuam como fortes regularizadores, reduzindo o sobreajuste (\textit{overfitting}) e viabilizando o treinamento de modelos complexos em domínios restritos \cite{perez2017effectiveness, shorten2019survey}.
%
% % data aug alem dos tradicionais (redes complexas) | como funciona modelagem em CN, comparacao com outros metodos e generica explicacao que CN gera descritores "robustos"
%
% Além das abordagens convencionais de transformação de imagem, a modelagem topológica oferece uma via promissora para o enriquecimento estrutural dos dados. Fundamentada na teoria dos grafos \cite{euler1736solutio}, a abstração de Redes Complexas mapeia a dinâmica de sistemas intrincados \cite{BOCCALETTI2006175}. Trabalhos pioneiros \cite{BACKES2013168} demonstraram que, ao modelar os pixels de uma textura como vértices e estabelecer arestas baseadas em um raio de distância euclidiana $r$, obtêm-se descritores topológicos robustos. A literatura evoluiu para incorporar o peso dessas arestas a partir de diferenças de intensidade \cite{ribas2020fusion}. Contudo, diferentemente das abordagens clássicas de análise de textura -- que comumente operam sobre imagens em escala de cinza, o presente projeto inova ao estender essa modelagem topológica para imagens RGB originais, capturando a complexa relação estrutural entre os canais de cor presentes em exames médicos.
%
% % explicacao sobre expansao do raio r e adaptacao para "entendimento" como serie temporal
% % rnn para capturar comportamento gerado pela CN
% A expansão progressiva do raio $r$ gera uma sequência ordenada de descritores, caracterizando uma dinâmica análoga a uma série temporal na qual as macroestruturas dependem da organização das microestruturas. Para capturar essa evolução espacial histórica, as Redes Neurais Recorrentes (RNNs) \cite{hochreiter1997long, cho2014learning} emergem como a solução ideal, interpretando a variação do grafo como passos sequenciais e gerando representações altamente discriminatórias.
%
% % two-stream networks como estratégia de alto impacto?? generico demais
% % fala sobre arquiteturas hibridas sucintamente
% % comenta sobre o late fusion rasamente tambem
% Para integrar essas dinâmicas topológicas com os atributos semânticos explícitos da imagem, arquiteturas híbridas de dupla via (\textit{Two-Stream Networks}) consolidaram-se como uma estratégia de alto impacto \cite{simonyan2014two}. A proposta deste trabalho é processar, em paralelo, as assinaturas estruturais (via RNN) e as características espaciais visuais da imagem (através de arquiteturas baseadas em convolução e atenção, inspiradas em modelos como a CoAtNet \cite{dai2021coatnet}). A união dessas vias por meio de fusão tardia (\textit{Late Fusion}) e concatenação de \textit{embeddings} visa estabelecer um classificador coeso e de precisão superior para cenários de dados limitados.
%

